\subsection{问题重述：一次往返与有条件修正}

题目给定的观测构成一条逐步加深的测厚链：四个附件按材料分成两组，分别对应同一块碳化硅晶圆和同一块硅晶圆。每块晶圆均有 $10^\circ$ 和 $15^\circ$ 入射条件下的观测，记录反射率随波数的变化。光学标定缺项统一列于第\ref{sec:optical-inputs}节。题目要求按一次返回条件建立外延层厚度关系，分别处理碳化硅与硅，并判断何时需要纳入多次往返影响。
% src:交接/数据档案.json:总览.附件总数;总览.独立晶圆数;总览.每块晶圆入射角度
% src:交接/数据档案.json:未随附件提供的建模输入

问题一从图\ref{fig:once-return}的表面反射和衬底界面首个返回出发，建立入射角、波数、材料光学性质与厚度的关系。本文将这两束光叠加后的光场称为两束截断场，即光进入外延层后只计入首次返回，后续往返留待问题三处理。题面给出的四组光谱是这一建模背景，问题一的模型由两束截断场构成，第三问再引入多次往返分母，区分各阶返回的影响。单次返回关系建立后，问题二再用同一片碳化硅晶圆的双角观测估计厚度。
% src:交接/题面契约.json:问题.0.需求条目.0.内容;硬约束清单.1.内容;硬约束清单.2.内容;硬约束清单.3.内容;硬约束清单.5.内容
% src:交接/图注素材_1.json:4.独有信息

同片碳化硅的坐标对应关系见第\ref{sec:common-coordinates}节；图\ref{fig:split}所对应的两角反射率均方根差为 2.14 个百分点，说明角度差异不能被一个共同振幅抹平。比较共享厚度与逐角厚度简单平均后，同片关系与折射率资料缺口促使模型共享厚度和随波数变化的折射率。这种安排还保留每个角度各自的光谱背景与条纹强度，使共同的厚度信息与角度造成的响应差异分别进入模型。
% src:交接/数据档案.json:跨文件关系.同片分组;跨文件关系.两两全量核验.0.公共键数;跨文件关系.两两全量核验.0.均方根差
% src:交接/实验记录.json:34.现象;34.决定

硅与碳化硅在问题三分成两条分支：硅晶圆无论是否满足多光束干涉条件，都必须给出厚度模型、算法和结果。碳化硅只有在确认多次往返出现且确实影响厚度精度时才修正。图\ref{fig:roundtrip}区分往返形成与可观测性，图\ref{fig:20}并列两种材料的答案。针对第\ref{sec:optical-inputs}节的观测缺项，谱形拟合与多光束可观测条件分项核对。三项任务的边界由此固定，下一节将用总体分析图串起数据入口、模型承接和结果出口。
% src:交接/题面契约.json:问题.1.解读;问题.2.解读;硬约束清单.13.内容;硬约束清单.14.内容
% src:交接/数据档案.json:未随附件提供的建模输入
% src:交接/图注素材_3.json:0.独有信息;5.独有信息
